\documentclass[12pt,a4paper]{article}
\usepackage[utf8]{inputenc}
\usepackage[T2A]{fontenc}
\usepackage[russian]{babel}
\usepackage{amsmath,amssymb,amsfonts}
\usepackage{geometry}
\geometry{left=2.5cm,right=2.5cm,top=2.5cm,bottom=2.5cm}
\usepackage{hyperref}
\usepackage{listings}
\usepackage{xcolor}
\usepackage{booktabs}
\usepackage{enumitem}

\lstset{
  language=Python,
  basicstyle=\ttfamily\small,
  keywordstyle=\color{blue},
  commentstyle=\color{gray},
  stringstyle=\color{red!70!black},
  showstringspaces=false,
  breaklines=true,
  frame=single,
  backgroundcolor=\color{gray!5},
}

\title{Code Review: Physical Optics в GOAD\\[0.3em]
\large\texttt{github.com/hballington12/goad}}
\author{Автоматический анализ}
\date{Март 2026}

\begin{document}
\maketitle
\tableofcontents
\newpage

% ====================================================================
\section{Обзор проекта}
% ====================================================================

GOAD (Geometric Optics with Aperture Diffraction) --- солвер рассеяния света
на крупных частицах ($x = ka \gg 1$), реализованный на Rust с Python-биндингами (PyO3).
Метод: трассировка лучей (geometric optics) + дифракция Фраунгофера на апертуре.

\textbf{Ключевые файлы с физикой:}
\begin{itemize}[nosep]
  \item \texttt{src/fresnel.rs} --- коэффициенты Френеля (отражение/пропускание)
  \item \texttt{src/snell.rs} --- закон Снеллиуса (комплексный показатель)
  \item \texttt{src/diff.rs}, \texttt{src/diff2.rs} --- дифракция Фраунгофера на апертуре
  \item \texttt{src/field.rs} --- амплитудная $\to$ матрица Мюллера, вращения поляризации
  \item \texttt{src/zones.rs} --- сечения (рассеяние, экстинкция, поглощение)
  \item \texttt{src/beam.rs} --- распространение лучей в поглощающих средах
  \item \texttt{src/orientation.rs} --- усреднение по ориентациям (Halton/Sobol)
  \item \texttt{src/convergence/} --- алгоритм Вэлфорда для адаптивной сходимости
\end{itemize}


% ====================================================================
\section{Корректные формулы}
% ====================================================================

% --------------------------------------------------------------------
\subsection{Фраунгоферов интеграл дифракции}
% --------------------------------------------------------------------

\textbf{Файлы:} \texttt{diff.rs:241--286}, \texttt{diff2.rs:594--637}

Дифракция на апертуре вычисляется как контурный интеграл по рёбрам апертуры
(теорема Грина для 2D):
\[
  F(\mathbf{k}) = \frac{k}{2\pi} \sum_{\text{edges}} e^{i \mathbf{k} \cdot \mathbf{r}_j}
  \left[ \alpha \frac{\cos\delta - \cos(\delta+\omega_1)}{?}
       - \beta  \frac{\cos\delta - \cos(\delta+\omega_2)}{?} \right]
\]
где $\delta = k_x x_j + k_y y_j$ --- фаза в вершине, $\alpha$, $\beta$ --- весовые коэффициенты рёбер.

\textbf{Статус:} \textcolor{green!50!black}{\textbf{Корректно.}}
Нормировка $k/(2\pi)$ правильная. Обработка краевых наклонов через \texttt{adjust\_mj\_nj()}.

% --------------------------------------------------------------------
\subsection{Закон Снеллиуса для комплексных показателей}
% --------------------------------------------------------------------

\textbf{Файл:} \texttt{snell.rs:55--91}

Реализация по Macke (1996), портированная из Fortran (\texttt{rt\_c.f90}):
\[
  \sin\theta_t = \frac{\sin\theta_i}{n_{\mathrm{rel}}^*}
\]
с учётом мнимой части показателя преломления через параметры:
\begin{align}
  k_1 &= \mathrm{Im}(m_1)/\mathrm{Re}(m_1), \quad
  k_2 = \mathrm{Im}(m_2)/\mathrm{Re}(m_2), \\
  n_{\mathrm{rel}} &= \frac{\mathrm{Re}(m_2)}{\mathrm{Re}(m_1)} \cdot
    \frac{1 + k_1 k_2}{1 + k_1^2}.
\end{align}

\textbf{Статус:} \textcolor{green!50!black}{\textbf{Корректно.}}
Проверенная реализация по рецензированному источнику. Есть тесты для нормального падения,
$30°$ и поглощающих сред.

% --------------------------------------------------------------------
\subsection{Оптическая теорема}
% --------------------------------------------------------------------

\textbf{Файл:} \texttt{zones.rs:349}

\[
  \sigma_{\mathrm{ext}} = \frac{4\pi}{k^2}\, \mathrm{Im}\bigl(S(\theta=0, \varphi=0)\bigr)
\]

\textbf{Статус:} \textcolor{green!50!black}{\textbf{Корректно.}}
Стандартная формула. Используется $k = 2\pi/\lambda$.

% --------------------------------------------------------------------
\subsection{Сечение рассеяния и параметр асимметрии}
% --------------------------------------------------------------------

\textbf{Файл:} \texttt{zones.rs:316--394}

\begin{align}
  \sigma_{\mathrm{sca}} &= \int_0^{2\pi}\!\!\int_0^{\pi}
    \frac{S_{11}(\theta,\varphi)}{k^2}\,\sin\theta\,d\theta\,d\varphi, \\
  g &= \frac{1}{\sigma_{\mathrm{sca}}}
    \int_0^{2\pi}\!\!\int_0^{\pi}
    \frac{S_{11}\cos\theta}{k^2}\,\sin\theta\,d\theta\,d\varphi, \\
  \omega &= \sigma_{\mathrm{sca}} / \sigma_{\mathrm{ext}}.
\end{align}

\textbf{Статус:} \textcolor{green!50!black}{\textbf{Корректно.}}

% --------------------------------------------------------------------
\subsection{Обратное рассеяние и деполяризация}
% --------------------------------------------------------------------

\textbf{Файл:} \texttt{zones.rs:366--387}

\begin{align}
  \sigma_{\mathrm{bs}} &= \frac{4\pi}{k^2}\, S_{11}(\theta = 180°), \\
  \delta_L &= \frac{S_{11}(180°) - S_{22}(180°)}{S_{11}(180°) + S_{22}(180°)}, \\
  \mathrm{LR} &= \sigma_{\mathrm{ext}} / \sigma_{\mathrm{bs}}.
\end{align}

\textbf{Статус:} \textcolor{green!50!black}{\textbf{Корректно.}}
Стандартные определения для лидарных приложений.

% --------------------------------------------------------------------
\subsection{Амплитудная матрица $\to$ матрица Мюллера}
% --------------------------------------------------------------------

\textbf{Файл:} \texttt{field.rs:76--157}

Преобразование $2\times 2$ комплексной S-матрицы в $4\times 4$ вещественную матрицу Мюллера:
\begin{align}
  M_{11} &= \tfrac{1}{2}(|S_{11}|^2 + |S_{12}|^2 + |S_{21}|^2 + |S_{22}|^2), \\
  M_{12} &= \tfrac{1}{2}(|S_{11}|^2 - |S_{12}|^2 + |S_{21}|^2 - |S_{22}|^2), \\
  M_{13} &= \mathrm{Re}(S_{11}S_{12}^* + S_{22}S_{21}^*), \\
  M_{14} &= \mathrm{Im}(S_{11}S_{12}^* - S_{22}S_{21}^*), \\
  &\;\vdots \notag \\
  M_{43} &= \mathrm{Im}(S_{22}S_{11}^* - S_{12}S_{21}^*), \\
  M_{44} &= \mathrm{Re}(S_{22}S_{11}^* - S_{12}S_{21}^*).
\end{align}

\textbf{Статус:} \textcolor{green!50!black}{\textbf{Корректно.}}
Все 16 элементов проверены. Знаки и множители $1/2$ правильные.

% --------------------------------------------------------------------
\subsection{Вращение базиса поляризации}
% --------------------------------------------------------------------

\textbf{Файл:} \texttt{field.rs:374--387}

Матрица поворота между базисами $(\hat{e}_\perp, \hat{e}_\parallel)$:
\[
  R = \frac{1}{\sqrt{|\det|}}
  \begin{pmatrix}
    \hat{e}_\perp^{\mathrm{out}} \cdot \hat{e}_\perp^{\mathrm{in}} &
    \hat{e}_\perp^{\mathrm{out}} \cdot \hat{e}_\parallel^{\mathrm{in}} \\
    -\hat{e}_\perp^{\mathrm{out}} \cdot \hat{e}_\parallel^{\mathrm{in}} &
    \hat{e}_\perp^{\mathrm{out}} \cdot \hat{e}_\perp^{\mathrm{in}}
  \end{pmatrix}
\]

\textbf{Статус:} \textcolor{green!50!black}{\textbf{Корректно.}}
Нормировка детерминантом обеспечивает унитарность.

% --------------------------------------------------------------------
\subsection{Поглощение в луче}
% --------------------------------------------------------------------

\textbf{Файл:} \texttt{beam.rs:199}

\[
  E \mathrel{*}= \exp\!\bigl(-2k\,\mathrm{Im}(n)\cdot d\bigr)
\]

\textbf{Статус:} \textcolor{green!50!black}{\textbf{Корректно.}}
Множитель 2 учитывает переход от интенсивности ($E^2$) к амплитуде.


% ====================================================================
\section{Обнаруженные проблемы}
% ====================================================================

% --------------------------------------------------------------------
\subsection{Проблема~1: Необъяснённая эмпирическая константа}
% --------------------------------------------------------------------

\textbf{Файл:} \texttt{diff.rs:84--98}

\textbf{Серьёзность:} \textcolor{red}{\textbf{Высокая}} (влияет на абсолютные сечения)

Масштабный коэффициент при переходе от геометрической оптики к дифракции:
\[
  \mathrm{scale} = \frac{\sqrt{\mathrm{CSA}}}{\sqrt{\Omega}} \cdot
  \underbrace{5.34464802915}_{\text{``bodge factor''}} \cdot
  \frac{\mathrm{phase}}{\lambda}
\]
где CSA~--- площадь поперечного сечения луча, $\Omega$~--- телесный угол.

Число $5.34464802915$ помечено в коде как <<empirical bodge factor>>
(подгоночный коэффициент). Оно не совпадает ни с одной стандартной константой:
\begin{itemize}[nosep]
  \item $2\sqrt{\pi} \approx 3.5449$
  \item $4\pi \approx 12.566$
  \item $\sqrt{4\pi} \approx 3.5449$
  \item $2\pi \approx 6.2832$
\end{itemize}

\textbf{Последствия:} Абсолютные значения сечений рассеяния могут быть
неправильно нормированы. Относительные величины (матрица Мюллера,
параметр асимметрии) не затрагиваются.

\textbf{Рекомендация:} Требуется аналитический вывод или ссылка на источник.

% --------------------------------------------------------------------
\subsection{Проблема~2: Коэффициенты Френеля для пропускания}
% --------------------------------------------------------------------

\textbf{Файл:} \texttt{fresnel.rs:36--47}

\textbf{Серьёзность:} \textcolor{orange}{\textbf{Средняя}} (требует проверки)

Реализация:
\begin{lstlisting}[language={}]
// p-polarization:
f11 = 2*n1*cos(th_i) / (n1*cos(th_t) + n2*cos(th_i))
// s-polarization:
f22 = 2*n1*cos(th_i) / (n1*cos(th_i) + n2*cos(th_t))
\end{lstlisting}

Стандартные формулы Френеля (Born \& Wolf, Hecht):
\begin{align}
  t_p &= \frac{2 n_1 \cos\theta_i}{n_2 \cos\theta_i + n_1 \cos\theta_t}, \\
  t_s &= \frac{2 n_1 \cos\theta_i}{n_1 \cos\theta_i + n_2 \cos\theta_t}.
\end{align}

Для p-поляризации знаменатель в коде: $n_1\cos\theta_t + n_2\cos\theta_i$,
а в учебнике: $n_2\cos\theta_i + n_1\cos\theta_t$.
Это \textbf{одно и то же} (коммутативность сложения), формула корректна.

Для s-поляризации: $n_1\cos\theta_i + n_2\cos\theta_t$ --- совпадает.

\textbf{Вывод:} При более тщательной проверке формулы оказались \textcolor{green!50!black}{\textbf{корректными}}.
Первоначальное подозрение снято.

% --------------------------------------------------------------------
\subsection{Проблема~3: Угол поля зрения в дифракции}
% --------------------------------------------------------------------

\textbf{Файл:} \texttt{diff2.rs:251}

\textbf{Серьёзность:} \textcolor{yellow!60!black}{\textbf{Низкая}} (фильтрация, не физика)

\begin{lstlisting}[language={}]
cos_fov = cos(fov_factor * 2.0 * 2.0 * PI / (k * aperture_dim))
\end{lstlisting}

Множитель $4\pi$ в аргументе косинуса не очевиден. Комментарий в коде
указывает на формулу \texttt{fov\_factor * $\lambda$/2 * r}, но реализация
имеет другую структуру.

\textbf{Последствия:} Влияет только на фильтрацию дифракционных вкладов
вне поля зрения. На финальный результат при правильном выборе \texttt{fov\_factor}
не влияет.

% --------------------------------------------------------------------
\subsection{Проблема~4: Матрица поляризации Карчевского}
% --------------------------------------------------------------------

\textbf{Файл:} \texttt{diff2.rs:337--378}

\textbf{Серьёзность:} \textcolor{yellow!60!black}{\textbf{Низкая}} (вероятно корректно, но сложно проверить)

Матрица преобразования между базисами поляризации падающего и рассеянного пучков:
\begin{lstlisting}[language={}]
frac = sqrt((1 - k_y^2) / (1 - K_y^2))
a1m = -K_z * frac       a1e = -k_z / frac
b2m = -k_z / frac       b2e = -K_z * frac
b1m = -k_x*k_y/frac + K_x*K_y*frac
a2e = -b1m
Matrix = [[0.5*(a1m+a1e), 0.5*b1m],
          [0.5*a2e, 0.5*(b2m+b2e)]]
\end{lstlisting}

Множитель $0.5$ --- усреднение TM и TE компонент. Формула основана на
Karczewski \& Wolf (1965), но без прямой ссылки в коде.

\textbf{Рекомендация:} Добавить ссылку на оригинальную статью.


% ====================================================================
\section{Сводная таблица}
% ====================================================================

\begin{center}
\begin{tabular}{llcl}
\toprule
\textbf{Формула} & \textbf{Файл} & \textbf{Статус} & \textbf{Примечание} \\
\midrule
Дифракция Фраунгофера & \texttt{diff.rs} & \textcolor{green!50!black}{\checkmark} & контурный интеграл \\
Закон Снеллиуса (компл.) & \texttt{snell.rs} & \textcolor{green!50!black}{\checkmark} & Macke 1996 \\
Оптическая теорема & \texttt{zones.rs} & \textcolor{green!50!black}{\checkmark} & $\sigma = 4\pi/k^2 \cdot \mathrm{Im}(S_0)$ \\
Сечения рассеяния & \texttt{zones.rs} & \textcolor{green!50!black}{\checkmark} & интеграл $S_{11}\sin\theta$ \\
Параметр асимметрии & \texttt{zones.rs} & \textcolor{green!50!black}{\checkmark} & $\langle\cos\theta\rangle$ \\
Деполяризация & \texttt{zones.rs} & \textcolor{green!50!black}{\checkmark} & $(S_{11}-S_{22})/(S_{11}+S_{22})$ \\
$S$-матрица $\to$ Мюллер & \texttt{field.rs} & \textcolor{green!50!black}{\checkmark} & все 16 элементов \\
Вращения поляризации & \texttt{field.rs} & \textcolor{green!50!black}{\checkmark} & унитарные \\
Поглощение в луче & \texttt{beam.rs} & \textcolor{green!50!black}{\checkmark} & $e^{-2k n'' d}$ \\
Френель (отражение) & \texttt{fresnel.rs} & \textcolor{green!50!black}{\checkmark} & $r_p$, $r_s$ \\
Френель (пропускание) & \texttt{fresnel.rs} & \textcolor{green!50!black}{\checkmark} & $t_p$, $t_s$ \\
\midrule
Масштаб дифракции & \texttt{diff.rs:87} & \textcolor{red}{\textbf{?}} & <<bodge factor>> 5.345 \\
Угол FOV & \texttt{diff2.rs:251} & \textcolor{yellow!60!black}{\textbf{?}} & множитель $4\pi$ \\
\bottomrule
\end{tabular}
\end{center}

% ====================================================================
\section{Общие выводы}
% ====================================================================

\begin{enumerate}
  \item \textbf{Физика в целом корректна.} Все стандартные формулы (Френель, Снеллиус,
        оптическая теорема, матрица Мюллера, сечения) реализованы правильно.

  \item \textbf{Главная проблема} --- необъяснённый подгоночный коэффициент 5.345 в
        масштабировании дифракционного вклада (\texttt{diff.rs:87}).
        Это влияет на абсолютные значения сечений, но не на относительные
        величины (матрица Мюллера, асимметрия).

  \item \textbf{Метод ограничен крупными частицами} ($x = ka \gg 1$), где
        геометрическая оптика применима. Для малых частиц ($ka \lesssim 10$)
        волновые эффекты (резонансы, дифракция) не описываются корректно ---
        здесь необходим BEM или DDA.

  \item \textbf{Хорошие практики:}
  \begin{itemize}[nosep]
    \item Проверенные алгоритмы из Fortran (Macke 1996)
    \item Тесты для критических функций (Снеллиус, Френель)
    \item Адаптивная сходимость (Welford + Halton/Sobol) --- мы заимствовали этот подход
  \end{itemize}

  \item \textbf{Конвенции:} $k = 2\pi/\lambda$, углы в радианах, предположительно $e^{+i\omega t}$.
\end{enumerate}

\end{document}
